% TOP_PROBLEM: {"id": 2306011, "title": "Research Problems in Function Theory — Problem 6.11", "category_id": 8, "category": {"id": 8, "name": "analysis", "display_name": "Analysis", "description": "Limits, continuity, calculus, and function theory.", "slug": "analysis", "order_index": 8, "created_at": "2026-07-31T15:26:25.671Z"}, "status": "open"}
% FIRST_POSTED: null
% ENTRY_KIND: statement_only
% Imported from: batch01.tex; UnsolvedMath ID 2306011
\documentclass[11pt]{book}
\usepackage{fontspec}
\setmainfont{Latin Modern Roman}
\setsansfont{Latin Modern Sans}
\usepackage{amsmath,amssymb,mathtools}
\usepackage{unicode-math}
\setmathfont{Latin Modern Math}
\usepackage{xurl}
\usepackage[hidelinks]{hyperref}
\newcommand{\sref}[2]{\hyperlink{source:#1:#2}{[S#2]}}
\newcommand{\cataloguescope}[1]{\paragraph{Mathematical question.}#1}
\begin{document}
% UnsolvedMath ID 2306011; source catalogue code AMR-022-6011
\setcounter{section}{2306010}
\section{Research Problems in Function Theory — Problem 6.11}\label{2306011}
{\small\sffamily UnsolvedMath ID 2306011 · Analysis}
\subsection{Definitions and mathematical statement}

\paragraph{Shared notation.}$\mathbb N=\{1,2,\ldots\}$, and $\mathbb Z,\mathbb Q,\mathbb R,\mathbb C$ denote the integers, rationals, reals and complex numbers. Any other conventions are specified locally.


Let $\mathbb D=\{z\in\mathbb C:|z|<1\}$, $\mathbb T=\partial\mathbb D$, and $\mathbb D^*=\{z\in\mathbb C:|z|>1\}$. Let $S$ be the class of analytic injective functions $f:\mathbb D\to\mathbb C$ normalized by $f(0)=0$, $f^{\prime}(0)=1$, with expansion $f(z)=z+\sum_{n=2}^{\infty}a_nz^n$ and $a_1=1$. A function in $S$ is starlike if its image is starlike about $0$: $tw\in f(\mathbb D)$ for every $w\in f(\mathbb D)$ and $0\le t\le1$. Call $f\in S$ convex when $f(\mathbb D)$ is a convex domain. For convex $f,g\in S$ and $0<\lambda<1$, put $h=\lambda f+(1-\lambda)g$.
\paragraph{Statement recorded in the supplied source.}
Is $h$ necessarily univalent and starlike about $0$? The original update reports a counterexample by Macgregor, together with the largest disk on which the starlikeness conclusion holds. \sref{2306011}{2}
\subsection{Short English statement}
Does every convex combination of two normalized convex disk mappings remain univalent and starlike? The source records a counterexample.
\subsection{Sources}
\begin{description}
\item[\hypertarget{source:2306011:1}{S1}] ULAM.AI and the UnsolvedMath Contributors, \emph{UnsolvedMath: A Curated Collection of Open Mathematics Problems}, record 2306011 (AMR-022-6011). CC BY 4.0. \url{https://huggingface.co/datasets/ulamai/UnsolvedMath}.
\item[\hypertarget{source:2306011:2}{S2}] Walter K. Hayman and Edward F. Lingham, \emph{Research Problems in Function Theory} (2018), arXiv:1809.07200. Problem 6.11, complete original question, all following updates, and the relevant chapter definitions. \url{https://arxiv.org/abs/1809.07200}.
\end{description}
\label{end:2306011}
\end{document}
